% !TEX root = ../main.tex
\section{Discussion}

Systematic reviews provide the evidentiary foundation for characterizing infectious disease epidemics, yet their production remains constrained by labor-intensive screening and downstream evidence extraction. In this study, we developed \metagent{}, a structured LLM-assisted workflow for systematic reviews of infectious disease modeling parameters. Across 22 COVID-19 and mpox source-review screening tasks comprising 12,758 candidate records and 668 ground-truth included studies, \metagent{} achieved an overall recall of 0.960 while reducing the overall \nns{} from 19.10 to 5.25. Performance was consistent across three epidemiological parameter domains, with recall reaching 0.984 for reproduction number, 0.935 for serial interval, and 0.957 for fatality. These findings suggest that LLM-assisted screening can substantially reduce the number of records requiring human review while preserving the sensitivity required in systematic review workflows.

The screening module was designed to reflect the eligibility logic used by domain experts. Rather than relying on a single binary relevance judgment, \metagent{} generates criterion-specific scores and rationales for disease, population, location, empirical-evidence, and parameter criteria before assigning an overall Strong, Possible, or Unlikely triage label. This structure reflects the fact that eligibility in infectious disease parameter reviews is rarely determined by one criterion alone. For example, an article may focus on the correct disease but report no empirical estimate of the target parameter, or it may estimate a relevant parameter using non-human or purely theoretical data. The intermediate relevance categories provide a transparent representation of the screening decision, supporting human oversight and reducing the risk of premature exclusion.

The framework was intentionally optimized for recall rather than precision. In systematic review screening, false negatives are generally more consequential than false positives, because missed eligible studies may bias downstream evidence synthesis and parameter estimation. \metagent{} therefore treats both Strong and Possible articles as relevant for manual review. This conservative strategy explains the moderate overall precision but is appropriate for a first-pass screening tool whose purpose is not to replace expert judgment, but to prioritize records and reduce the screening burden. This is particularly critical in resource-constrained emergency settings, such as guiding evidence-based public health responses during sudden outbreaks. The reduction in \nns{} indicates that reviewers would need to screen fewer records for each eligible study retained, allowing reviewers to focus on a smaller and more enriched set of candidate studies.

Performance varied across parameter domains, reflecting differences in how epidemiological concepts are reported in the literature. Reproduction-number reviews achieved the highest recall (0.984), with fatality (0.957) and serial-interval (0.935) reviews also retaining high sensitivity. Although serial-interval recall was the lowest among the three domains, inspection of false negatives indicated that some eligible studies reported the target parameter only in full-text tables, appendices, or secondary analyses rather than in the abstract. This limitation is inherent to title-and-abstract screening and is not unique to LLM-based methods. Reproduction number tasks showed high recall but higher \nns{} than serial interval tasks, likely because broad mentions of transmissibility and epidemic growth make it more difficult to distinguish articles that estimate reproduction numbers from those that merely discuss them. These findings highlight that screening difficulty is partly determined by the semantic specificity of the target parameter. An analogous parameter-specific gradient was observed in full-text coding: fatality estimates required additional alignment of denominators and population strata before they could be compared with the source-review references, whereas serial-interval and reproduction-number estimates could be coded against more uniformly defined estimands. The screening- and coding-side challenges therefore share a common origin in the heterogeneity of how each target parameter is defined and reported, rather than in the discriminative capacity of the underlying model.

\metagent{} also addresses an important limitation of generic LLM applications in systematic review automation. Previous fully automated GPT-based approaches have shown limited ability to identify eligible articles in systematic reviews, underscoring the need for task-specific workflow design. In contrast, \metagent{} does not require domain-specific model retraining, but improves screening performance through structured prompting and explicit eligibility criteria applied to a uniform title-and-abstract input. This suggests that carefully designed LLM workflows, augmented with epidemiological domain knowledge and review-specific eligibility criteria, may provide practical gains even without model fine-tuning. Such a design may be especially valuable in rapidly evolving outbreak contexts, where timely evidence synthesis is needed but labeled training data may be limited.

Beyond screening, the full-text parameter coding and extraction experiment indicates that LLM-assisted workflows can support later stages of infectious disease meta-analysis when the included studies are fixed. The coding component did not attempt to replace full prospective review conduct; instead, it tested whether target-parameter evidence could be localized, coded, extracted, and aggregated to align with estimates reported by the source reviews. This distinction is important because parameter coding and extraction depend not only on identifying relevant papers, but also on aligning denominators, populations, estimands, and synthesis rules. The close agreement with source-review estimates therefore supports the feasibility of codebook-guided parameter coding and extraction, while preserving the need for expert adjudication when review definitions are ambiguous.

\metagent{} provides valuable assistance to medical experts and systematic reviewers by streamlining the literature mining process and offering superior temporal efficiency compared to purely manual efforts. However, the feasibility of AI integration remains constrained by several unresolved concerns, including algorithmic opacity, compromised reproducibility across temporally distinct queries, and limited capacity for context-dependent interpretation.\cite{ref28} Thus, while AI offers considerable potential as an adjunctive instrument for accelerating review workflows, a critical examination of its methodological boundaries and the enduring necessity of human oversight warrants further investigation.\cite{ref29,ref30}

Despite these strengths, several limitations should be considered. First, the ground truth was derived from the inclusion lists of published systematic reviews rather than from a newly adjudicated gold standard. Published reviews may contain retrieval omissions, apply eligibility criteria that are incompletely reported, or include studies identified through databases and supplementary searches not reproduced in our benchmark. As a result, recall and precision estimates should be interpreted as performance against review-derived ground truth, not as definitive measures against an exhaustive corpus of all eligible studies. Second, our benchmark was constructed primarily from PubMed, whereas the original systematic reviews often searched multiple bibliographic databases. This structural asymmetry may limit the generalizability of the reported performance to complete real-world systematic review workflows involving Embase, Web of Science, Scopus, preprint servers, and manual citation tracking. Third, the full-text screening comparison and parameter coding analysis were evaluated as controlled experiments rather than as a fully prospective end-to-end review. Articles labeled Strong or Possible still require expert review, and final inclusion decisions remain the responsibility of human reviewers. Fourth, the method depends on the information available in titles, abstracts, accessible full texts, and parsed tables. Studies that report relevant parameters only in supplementary materials, images, or poorly parsed tables may still be missed or incompletely coded. The focused full-text screening comparison suggests that access to full-text evidence can improve both sensitivity and specificity when relevant evidence is absent from the abstract, but this benefit was conditional on successful full-text retrieval and parsing. In practice, many records lack informative abstracts, are not available through PMC Open Access or other accessible full-text sources, or contain the relevant parameter only in tables, appendices, or sections that are difficult to localize automatically. These constraints point to an infrastructure-level limitation of current evidence synthesis workflows: more open, structured, and machine-readable repositories of biomedical full text, tables, supplementary materials, and review-level extraction targets would likely improve both LLM-assisted screening and downstream parameter coding. Incorporating broader full-text retrieval and multimodal table understanding may reduce this source of error, but would introduce additional challenges related to access, document parsing, and computational cost.

More broadly, integrating LLMs into evidence synthesis raises methodological and operational concerns. LLM outputs may vary across model versions, prompts, and time, affecting reproducibility. The internal reasoning process of proprietary models may not be fully transparent, and model behavior may shift as systems are updated. These issues are particularly important for systematic reviews, where reproducibility and auditability are central methodological principles. The structured decision design of \metagent{} improves interpretability at the decision level, but does not eliminate the need for careful documentation, version control, prompt archiving, and human verification.

Future work should evaluate \metagent{} prospectively in live systematic review projects and across a wider range of pathogens, study designs, and epidemiological parameters. Multi-database benchmarks would better reflect real-world search workflows, while prospective validation could assess how the system performs when screening newly published literature not available during model pretraining. Further development should also explore integration with full-text screening, citation chasing, uncertainty calibration, and active-learning prioritization. Combining LLM-based semantic assessment with conventional information retrieval methods may further improve both sensitivity and efficiency. In parallel, the coding and extraction component warrants dedicated extension. Promising directions include multimodal table and figure understanding for parameter values that are not represented in running text, automatic harmonization of denominators and estimands across reviews, transferable codebooks that generalize to new pathogens or parameter families, and uncertainty-aware aggregation that explicitly propagates extraction confidence into downstream meta-analytic estimates. Joint optimization of retrieval and parameter-aware coding within a single human-in-the-loop framework may further narrow the gap between LLM-assisted output and review-ready evidence synthesis.

In summary, \metagent{} demonstrates that structured LLM-assisted workflows can support systematic reviews of infectious disease modeling parameters. The framework preserves high screening recall, lowers \nns{}, and assists full-text parameter coding and extraction as a modular component of evidence synthesis. These results support the use of structured, human-in-the-loop LLM systems as complementary infrastructure for systematic reviews, particularly in public health settings where timely and reliable parameter evidence is essential.
